\corrigesChapitre{10}{Exploiter la gravitation universelle}
\par\noindent\begin{minipage}{\linewidth}
\begin{corrige}{1}{Nature de l’action}
\begin{questions}
\item Non : elle agit à distance.
\item Oui.
\end{questions}
\end{corrige}
\end{minipage}
\par\noindent\begin{minipage}{\linewidth}
\begin{corrige}{2}{Deux objets}
\begin{questions}
\item Oui, avec une force de même valeur.
\item Sur la Terre.
\end{questions}
\pcFigureEx{figures/rev-cor-gravitation.pdf}{Deux flèches opposées, appliquées à deux corps différents ; elles ne se compensent pas sur le seul corps A.}
\end{corrige}
\end{minipage}
\par\noindent\begin{minipage}{\linewidth}
\begin{corrige}{3}{Unités}
\begin{questions}
\item kg et m.
\item $6{,}4\times10^6\unite m$.
\end{questions}
\end{corrige}
\end{minipage}
\par\noindent\begin{minipage}{\linewidth}
\begin{corrige}{4}{Quelle distance ?}
\begin{questions}
\item La distance entre leurs centres dans le modèle.
\item Il faut aussi connaître les rayons pour retrouver cette distance centrale.
\end{questions}
\pcFigureEx{figures/rev-cor-centres.pdf}{La distance d de la loi gravitationnelle est mesurée entre les centres, pas entre les surfaces.}
\end{corrige}
\end{minipage}
\par\noindent\begin{minipage}{\linewidth}
\begin{corrige}{5}{Une masse doublée}
\begin{questions}
\item F double.
\item F est multipliée par quatre.
\end{questions}
\end{corrige}
\end{minipage}
\par\noindent\begin{minipage}{\linewidth}
\begin{corrige}{6}{Une distance doublée}
\begin{questions}
\item F est divisée par quatre.
\item F est multipliée par quatre.
\end{questions}
\pcFigureEx{figures/rev-cor-inverse-carre.pdf}{Le facteur de distance est mis au carré au dénominateur ; à trois fois la distance, la force est divisée par neuf.}
\end{corrige}
\end{minipage}
\par\noindent\begin{minipage}{\linewidth}
\begin{corrige}{7}{Une force modèle}
\begin{questions}
\item $6{,}67\times10^{-11}\unite N$.
\item Non : elle est très faible, même si un appareil scolaire ne la détecte pas.
\end{questions}
\end{corrige}
\end{minipage}
\par\noindent\begin{minipage}{\linewidth}
\begin{corrige}{8}{Une comparaison rapide}
\begin{questions}
\item $36/9=4\unite N$.
\item $36/4=9\unite N$.
\end{questions}
\pcFigureEx{figures/rev-cor-inverse-carre.pdf}{Le facteur de distance est mis au carré au dénominateur ; à trois fois la distance, la force est divisée par neuf.}
\end{corrige}
\end{minipage}
\par\noindent\begin{minipage}{\linewidth}
\begin{corrige}{9}{Depuis l’orbite}
\begin{questions}
\item $d=6\,800\unite{km}=6{,}8\times10^6\unite m$.
\item Toute la distance en mètres, donc $(6{,}8\times10^6)^2$.
\end{questions}
\par\smallskip\noindent\textbf{Pour comprendre.}\enspace $d=R+h=6\,400+400=6\,800\unite{km}=6{,}8\times10^6\unite m$. Employer seulement l’altitude ferait partir la mesure de la surface et non du centre.\par
\pcFigureEx{figures/rev-cor-altitude.pdf}{Additionner rayon et altitude avant de convertir en mètres ; le point de départ est le centre de la Terre.}
\end{corrige}
\end{minipage}
\par\noindent\begin{minipage}{\linewidth}
\begin{corrige}{10}{Un petit calcul}
\begin{questions}
\item $F=3{,}335\times10^{-9}\unite N$.
\item Elle double : $6{,}67\times10^{-9}\unite N$.
\end{questions}
\end{corrige}
\end{minipage}
\par\noindent\begin{minipage}{\linewidth}
\begin{corrige}{11}{Trois facteurs}
\begin{questions}
\item Par $2/4=1/2$.
\item La masse au numérateur apporte 2 ; le carré de la distance au dénominateur apporte 4.
\end{questions}
\par\smallskip\noindent\textbf{Pour comprendre.}\enspace Le facteur de masse est 2, celui du carré de la distance est $2^2=4$ ; le facteur total est donc $2/4=0{,}5$. Traiter séparément les modifications avant de les combiner.\par
\end{corrige}
\end{minipage}
\par\noindent\begin{minipage}{\linewidth}
\begin{corrige}{12}{Un faux équilibre}
\begin{questions}
\item Elles s’exercent sur deux objets distincts et ne se compensent pas sur un seul.
\item L’attraction exercée par la Terre sur la Lune, avec les autres actions éventuelles du modèle.
\end{questions}
\pcFigureEx{figures/rev-cor-gravitation.pdf}{Deux flèches opposées, appliquées à deux corps différents ; elles ne se compensent pas sur le seul corps A.}
\end{corrige}
\end{minipage}
\par\noindent\begin{minipage}{\linewidth}
\begin{corrige}{13}{Impesanteur}
\begin{questions}
\item La station et les occupants restent attirés par la Terre.
\item Ils sont en chute commune autour de la Terre, hors forces de contact de soutien habituelles.
\end{questions}
\pcFigureEx{figures/rev-cor-orbite.pdf}{Le satellite et son contenu sont en chute commune ; l’impesanteur ne signifie pas absence de gravitation.}
\end{corrige}
\end{minipage}
\par\noindent\begin{minipage}{\linewidth}
\begin{corrige}{14}{Retrouver g}
\begin{questions}
\item $g=588/60=9{,}8\unite{N/kg}$.
\item $m=98/9{,}8=10\unite{kg}$.
\end{questions}
\end{corrige}
\end{minipage}
\par\noindent\begin{minipage}{\linewidth}
\begin{corrige}{15}{La précision du modèle}
\begin{questions}
\item La précision numérique affichée ne corrige pas l’approximation des données.
\item Les corps sont assimilés à des sphères ou des masses ponctuelles pour employer la distance entre centres.
\end{questions}
\end{corrige}
\end{minipage}
\clearpage
\par\noindent\begin{minipage}{\linewidth}
\phantomsection\label{cor-defi-10}\pcSolutionDefi{10}{Autour d’une planète}\begin{questions}
\item $F\approx4{,}08\times10^3\unite N$, soit environ $4{,}1\times10^3\unite N$.
\item Environ $1{,}02\times10^3\unite N$, par division par quatre.
\item Les valeurs sont égales et les sens opposés, mais les forces s’appliquent à deux masses très différentes. Égalité de force ne signifie pas évolution de mouvement identique.
\end{questions}
\par\smallskip\noindent\textbf{Pour comprendre.}\enspace $F=6{,}67\times10^{-11}\times(6{,}0\times10^{24})\times500/(7{,}0\times10^6)^2\approx4{,}1\times10^3\unite N$. Doubler la distance divise par quatre : environ $1{,}0\times10^3\unite N$.\par
\pcFigureEx{figures/rev-cor-satellite.pdf}{Distances entre centres et masses en unités SI ; la force diminue lorsque la distance augmente.}
\end{minipage}
\par\noindent\begin{minipage}{\linewidth}
\phantomsection\label{cor-activite-10}\pcSousTitre{Pistes pour l’investigation}
Pour d, 2d et 3d, les forces sont dans les rapports 1, 1/4 et 1/9. Si une seule masse double à distance fixée, F double. Le graphe F(d) est courbe ; un graphe de F en fonction de 1/d² serait linéaire pour des masses fixées. Le tableur vérifie les conséquences du modèle, il ne constitue pas une nouvelle preuve expérimentale de la loi.\par\end{minipage}
